% Page 071

\seite{71}

in der Bürgermeistereiratswahl gewählt wurde.\ob
Die großen Vorbereitungen – wie grundsätzlich – gab es für\ob
bei die Gemeinderatswahlen. Zunächst wurden einmal\ob
5 Wahlvorschläge aufgestellt, nämlich Wahlvorschlagsliste:\ob
1.\ Hoffmann,  2.\ Schneider,  3.\ Kasbach,  4.\ Weinandy-Schütz

\pend
\abschnitt{5 Stich – Wahlen –}
\pstart

die drei Letzteren wollten einen Wahlvorschlagsliste zu\ohb
rückgewesen sein, mit der Begründung, daß bei zu spät\ohb
einen Vorschreiben des Wahlausschalters eingeleiert wor\ohb
den. – Eingelaufen sind die Wahlvorschläge 3 u 4 nach\ob
der Uhr der Benachrichtigung genau um 5.05 Uhr – wie\ohb
hing dem Bürgermeister- Wilk. und Parola ist der Uhr\ohb
abgleich daran. Der Wahlvorschlag 5 ist 6.05 Uhr – also\ob
wirklich zu spät – am 17 Tage vor der Wahl abgege\ohb
ben worden. Da es wirklich begründet worden war\ohb
name, daß die Wahlvorschläge 3 u 4 pünktlich eingelau\ohb
fen worden, konnte in der der Wahlhandlung D\ohb
Grube formell gefragt haben die Wahlvorschläge angenommen\ob
worden sind, die Wahl unten\unsicher{} ist durch dieses Verfahren\ob
wann würdig in auf die Wahlvorschläge Kasbach\ob
und Weinandy-Schütz geschädigt, daß der Gemeinderats\ohb
wahl entgültig als Stimmen abgegeben wurden\ob
117.\ Davon entfielen auf die Wahlvorschläge:\ob
1.\ Hoffmann: 23 Stimmen, 2.\ Schneider 24 Stimmen,\ob
3.\ Kasbach:  37 Stimmen, 4.\ Weinandy-Schütz:\ob
18 Nummern. Ungültige Stimmen waren\ob
5 vorhanden, woran die Leisten absichtlich ungül\ohb
tig gemacht worden. Ein Stimmzettel enthält die\ob
Formen von je 2 Hauptspalten, Parteien und Kandidaten\ohb
plan. Dieser Stimmzettel muß als eine Umschlag\ohb
mappe der öffentlichen Verfassung umgegeben wor\ohb
den, es waren wenig Streit, jeder nahm seine\ob
Freiheit, so ein geheimes Wahl so nicht der offiziel\ohb
len sein dürften. Im Kreisteil hin die Umschlag\ohb
pflicht am anderen aufzufällen. Vieler Menschen\unsicher\ob
fiel sich unfreundlich über die Möglichkeit eines geheimen
\pend
